\section{MATERIALS AND METHODS}\label{materials-and-methods}

\subsection{Data}\label{data}

The headline evaluation matrix runs entirely against
Synthea-generated FHIR Bundles \autocite{Walonoski2018}, a freely-
redistributable synthetic patient dataset. Cohort generation is
fully reproducible from a pinned Synthea version and seed (see
\emph{Cohort and label space} below).

\subsection{Cohort and label space}\label{cohort-and-label-space}

We restrict the cohort to ICD-10-CM codes in the CMS ACCESS Model
FHIR Implementation Guide v0.9.6 \autocite{CMS2026}. The Implementation
Guide defines two disease groups: \textbf{CKM} (cardiometabolic, covering
diabetes, atherosclerotic cardiovascular disease, and CKD stage 3)
and \textbf{eCKM} (extended cardiometabolic, covering hypertension,
dyslipidemia, prediabetes, and obesity). The headline matrix is
evaluated on a single Synthea-generated cohort of \textbf{125 unique FHIR
Conditions}, each materialized into four degradation modes for a
total of 500 EvalRecord items. Cohort generation is fully
deterministic: Synthea v4.0.0 \autocite{Walonoski2018} is run at a pinned
commit SHA with \texttt{seed=42}, producing patient bundles whose
Conditions are extracted, deduplicated per (patient, ICD code), and
filtered to ACCESS-scope codes via bundled value-set definitions.
Twelve unique ICD-10-CM codes appear in the cohort: obesity (\textasciitilde24\%),
prediabetes (\textasciitilde20\%), type 2 diabetes variants (\textasciitilde31\%), essential
hypertension (\textasciitilde14\%), and lipid disorders (\textasciitilde10\%).

\subsection{Models evaluated}\label{models-evaluated}

The headline matrix comprises 28 model configurations. These are:
three string-matching baselines (exact, fuzzy, TF-IDF), one frozen
sentence-transformer retrieval baseline, and 24 frontier LLM
configurations across three providers and three prompting modes.
Each LLM appears in \texttt{zeroshot}, \texttt{constrained}, and \texttt{rag}
configurations.

\begin{itemize}
\tightlist
\item
  \textbf{Anthropic} (3 × 3): Claude Haiku 4.5, Sonnet 4.6, Opus 4.7
\item
  \textbf{OpenAI} (2 × 3): GPT-5.5 (\texttt{gpt-5.5-2026-04-23}) and GPT-4o-mini
  \autocite{OpenAIModels2026}
\item
  \textbf{Google} (3 × 3): Gemini 2.5 Pro, Gemini 2.5 Flash, and Gemini 3
  Flash Preview (preview-status, disclosed as such) \autocite{GeminiModels2026}
\end{itemize}

Gemini 3.1 Pro Preview was excluded due to insufficient daily quota
at our API tier; a partial-coverage analysis appears in
Supplementary §S2. Per-token pricing for cost calculations was
sourced from each provider's documentation as of the evaluation date
\autocite{OpenAIPricing2026,GeminiPricing2026}.

\subsection{Prompting modes}\label{prompting-modes}

All LLMs are evaluated under three prompting modes that vary in how
much external structure is provided.

\begin{itemize}
\tightlist
\item
  \textbf{\texttt{zeroshot}}: a generic system prompt instructs the model to
  return the most likely ICD-10-CM code for the input, with no
  candidate list.
\item
  \textbf{\texttt{constrained}}: the system prompt additionally includes the
  full list of CMS ACCESS-scope candidate codes (with display
  strings) and instructs the model to choose only from that list.
  Mechanical schema enforcement (Anthropic forced tool use, OpenAI
  structured outputs, Google \texttt{response\_schema}) is applied across
  all modes; in \texttt{constrained} mode the menu itself is also passed
  in-prompt.
\item
  \textbf{\texttt{rag}}: a frozen sentence-transformer (\texttt{all-MiniLM-L6-v2})
  retrieves the top-k = 10 ACCESS-scope candidate codes most similar
  to the per-record input and injects only those into a
  constrained-style prompt.
\end{itemize}

The verbatim system and user prompts and per-provider tool-use and
schema configurations are reproduced in Supplementary §S1.

\subsection{Outcome taxonomy}\label{outcome-taxonomy}

Each per-record top-1 prediction is classified into one of six
mutually exclusive, exhaustive outcome buckets, ordered from best
to worst by deployment relevance:

\begin{enumerate}
\def\labelenumi{\arabic{enumi}.}
\tightlist
\item
  \textbf{\texttt{exact\_match}}: predicted code equals ground truth.
\item
  \textbf{\texttt{category\_match}}: predicted code is in the same ICD-10-CM
  3-character category as ground truth (e.g., E11.0 vs.~E11.9, both
  Type 2 diabetes).
\item
  \textbf{\texttt{chapter\_match}}: same ICD-10-CM chapter, different category.
\item
  \textbf{\texttt{out\_of\_domain}}: real ICD-10-CM code with no hierarchical
  relation to ground truth.
\item
  \textbf{\texttt{no\_prediction}}: model returned no usable prediction (empty
  array, refusal, or transient API failure that exhausted the SDK
  retry budget). The model fabricated nothing; it abstained.
\item
  \textbf{\texttt{hallucination}}: predicted code does not exist in the FY2026
  CMS ICD-10-CM tabular list (mechanically validated against the
  bundled validator).
\end{enumerate}

The \texttt{no\_prediction} and \texttt{hallucination} buckets capture
qualitatively different failure modes and are reported separately.
From a downstream-error perspective, abstention is generally
preferable to fabrication, since an abstaining model emits no
spurious code that downstream systems silently mishandle. Earlier work \autocite{Soroush2024}
collapses these into a single error bucket; the six-way split
surfaces the distinction that conflation hides.

\subsection{Cost tracking}\label{cost-tracking}

For every model invocation in our matrix the per-prediction CSV
records token usage, latency, and \texttt{cost\_usd} computed at runtime
against the versioned pricing snapshot in \texttt{configs/pricing.yaml}
\autocite{OpenAIPricing2026,GeminiPricing2026}. These infrastructure
metrics are recorded even when the prediction is wrong and serve
as the basis for the cost-per-correct normalization reported in
Results. Cost-per-correct (USD per exact-match outcome) combines per-call
price and per-call accuracy into a single number; we report it
alongside per-call cost because the two diverge when accuracy
drops.

\subsection{Statistical analysis}\label{statistical-analysis}

Per-(model, mode) outcome rates are reported as point estimates
with 95\% Wilson confidence intervals, which we prefer over
normal-approximation intervals because both small N and rates near
zero or one are common in our matrix (per-cell N = 125, and
constrained-mode hallucination rates approach 0\%). Within-model
paired comparisons across zero-shot, constrained, and RAG modes on
identical inputs are reported as exact McNemar tests on discordant
prediction pairs. Across-model and across-mode comparisons are
reported descriptively with overlapping or non-overlapping CIs. We
deliberately do not perform null-hypothesis significance tests
across the full grid, because the intended interpretation is
comparative deployment relevance rather than detection of any
single effect.
